\section{The coefficient of the signed error}\label{sec:principal}
In the formula for the limit, $\mathcal C_0$ collects the adjustment on the initial board and the sum of
the insertion defects up to the stopping time. The term $P_0-S$ measures the error remaining among the
squares available at that time.

Let $x$ be the fraction of rows still empty. At $x=1$ the initial board is empty, and as $x$ approaches
zero only the last rows remain. The function $d(x)$ sums sixteen contributions to the mean second-order
defect of an insertion in the predictor~\cite{pc}. Appendix~\ref{app:interior} identifies these
contributions. The line prefactor $\mathcal V(x)$ and the rook correction $F(x)$ are defined in
Appendix~\ref{app:predictor}. The following two statements are proved \status{D}, given the results
imported in Section~\ref{sec:importados} and with two computer-assisted steps \status{O}; the labels
are explained in Section~\ref{sec:claim-ledger}. Here and below, the \emph{physical profile} is the
density $\nu_1$ of a complete solution with its diagonal vacancies $u_f$
(Section~\ref{sec:profile-data}); more generally, \emph{physical} marks an object as it occurs on the
board, as opposed to the approximants, continuations, and auxiliary grids used in the proof. The
\emph{ordinary board} is the usual $n\times n$ board, empty at the start, as opposed to the residual
boards of a history (Section~\ref{sec:importados}) and to the rook masks (Section~\ref{sec:torres});
this is the usage of~\cite{pc,tc}.

\begin{teorema}[Main theorem]\label{teo:principal}
For the ordinary board the limit
\begin{align}
 \cerr&=\lim_{n\to\infty}nR_Q(n)=\mathcal C_0+P_0-S,\label{eq:main}\\
 \mathcal C_0&=F(1)-\tfrac1{12}+\int_0^1 \frac{\numerr(x)}{x^2}\,dx,\notag\\
 \numerr(x)&=d(x)+G(x)-xG'(x),\qquad G(x)=\mathcal V(x)+F(x)-\tfrac1{12},\notag\\
 P_0&=P_T+J,\qquad J=-B_2+L_2,\notag
\end{align}
exists, and the integral in~\eqref{eq:main} converges absolutely.
\end{teorema}
On the full mask the scaled rook kernel is constant, so \(F(1)=0\) (Section~\ref{sec:profile-data}). The
terms are defined as follows: $d$ in~\eqref{eq:d-suma}; $F$ and $\mathcal V$ in~\eqref{eq:F-def}
and~\eqref{eq:V-def}; $P_T=\lim_{b\downarrow0}P_T(b)$ with $P_T(b)$ as
in~\eqref{eq:pt-frontera-fija}; $B_2$ and $L_2$ in~\eqref{eq:J-def}, through the state functionals
of~\eqref{eq:estado-def}; and $S$ in~\eqref{eq:S-def}. Here
$P_T$ is the rook contribution at arrival. To obtain $B_2$ and $L_2$ from the weight path of a collinear
pair, we sum over both orientations. Within each orientation we weight the four availability states of
the other two vertices. $S$ accounts for the interaction of the two pairs in a complete square.

\begin{corolario}\label{cor:derivada}
For the physical profile, $\numerr(0)=0$ and $\numerr'(0)=d'(0)=0$. Moreover,
\[
 \log Q(n)=\log Q_{\rm pred}(n)+\cerr/n+o(n^{-1}),\qquad
 Q(n)/Q_{\rm pred}(n)=1+\cerr/n+o(n^{-1}).
\]
\end{corolario}
Multiplying the predictor by \(\exp(\cerr/n)\) leaves a relative error of \(o(1/n)\). The corollary gives
no threshold beyond which the correction improves the result at every size. The proof takes the truncated
integral at a fixed cut-off $b>0$. Comparing two cut-offs gives its limit, and this convergence yields
$\numerr(0)=0$ and $\numerr'(0)=0$ (Appendix~\ref{app:closing}). Conjecture~10.1 of~\cite{pc} states that
there exist $C$ and $n_0$ such that $|\log Q(n)-\log Q_{\rm pred}(n)|\le C/n$ for all $n\ge n_0$;
Lemma~\ref{lem:previa} proves it. The corollary identifies the signed coefficient at that rate. The
proof provides an expression for the coefficient; certifying the sign suggested by its evaluation still
requires a sufficiently narrow interval.

\paragraph{Reading of the integral away from the endpoint.} Lemma~\ref{lem:regularidad-extremo} gives the
analyticity of $G$ only near $[0,x_0]$. On $[x_0,1]$ the proof uses the summand $(G-xG')/x^2$ only
through its antiderivative $-G(x)/x$: the interior identity of Appendix~\ref{app:interior} is
$\int_b^1d(x)x^{-2}\,dx+G(b)/b-G(1)$. Accordingly, the part of the integral in~\eqref{eq:main} over
$[x_0,1]$ is understood as $\int_{x_0}^1d(x)x^{-2}\,dx+G(x_0)/x_0-G(1)$. This requires only that $d$ be
continuous there (Appendix~\ref{app:interior}) and agrees with the classical integral whenever $G$ is
continuously differentiable on $[x_0,1]$. The same reading applies to the split in
Section~\ref{sec:presupuesto}.

Section~\ref{sec:signedmap} explains the signs in the formula. The table in
Section~\ref{sec:claim-ledger} identifies the sources for each step and the scope of its checks.
